\documentclass[10pt,a4paper]{amsart}
\usepackage[margin=19mm]{geometry}
\usepackage{amsmath,amssymb,mathtools}
\usepackage[scr=rsfs,cal=euler]{mathalfa}
\usepackage{xfrac}
\usepackage[unicode,hidelinks,bookmarks=false]{hyperref}
\newcommand{\ie}{i.e.\ }
\newcommand{\cf}{cf.\ }
\newcommand{\resp}{respectively\ }
\newcommand{\Subsection}{\subsection}
\providecommand{\nobreakdash}{\nobreak\hbox{-}}
\setlength{\emergencystretch}{2em}
\multlinegap0pt
\usepackage{booktabs}
\usepackage{array}
\usepackage{url}
\usepackage{multirow}
\usepackage{tikz}
\usetikzlibrary{arrows.meta, positioning}
\usepackage{subfig}
\usepackage{enumitem}
\definecolor{copenblue}{RGB}{0,102,204}
\definecolor{copenlight}{RGB}{173,216,230}
\definecolor{copendark}{RGB}{0,51,102}
\newtheorem{theorem}{Theorem}
\newtheorem{lemma}{Lemma}
\newtheorem{corollary}{Corollary}
\newtheorem{definition}{Definition}
\newtheorem{proposition}{Proposition}
\newtheorem{assumption}{Assumption}
\newtheorem{remark}{Remark}
\newcommand{\E}{\mathbb{E}}
\newcommand{\Prob}{\mathbb{P}}
\newcommand{\R}{\mathbb{R}}
\begin{document}
\title[Structural Decomposability of Encrypted Traffic Side-Channel]{Structural Decomposability of Encrypted Traffic Side-Channel Leakage}
\author{Guangjie Liu, Guang Cheng, Weiwei Liu}
\date{}
\maketitle
\noindent\emph{Source and licence.} Structural Decomposability of Encrypted Traffic Side-Channel Leakage. Version arXiv:2609.19036v1 (CC BY 4.0). This material is distributed under the Creative Commons Attribution 4.0 International licence, http://creativecommons.org/licenses/by/4.0/, as recorded in the arXiv OAI licence metadata for this exact version and shown on the abstract page https://arxiv.org/abs/2609.19036v1. Full rights notice: licenses/arxiv-2609.19036-CC-BY-4.0.txt.\par\smallskip
\noindent\textit{Editorial scope.} Counted unit: the original \texttt{proposition} environment \texttt{prop:fisher-chain} at source lines 895--912 together with its complete proof at lines 914--928; the delivered document keeps source lines 742--929 so that every referenced label and every citation resolves inside it. Native editable source is the paper's own concatenated TeX; the AMS wrapper is editorial formatting only. No publisher class, upstream script or unrelated artwork is redistributed.
\section{Information Geometry Framework}\label{sec:ig}
% ============================================================

\textbf{This section answers a purely side-channel question}: when the
defense knob is adjusted by a small step, how does the attacker's
separability change?
The core of the answer is that the KL divergence (the measure of
distributional separability) has a quadratic local approximation governed
by the Fisher information matrix.
Therefore, the size of the Fisher matrix eigenvalues in the defense
parameter direction directly characterizes ``how much separability change
the same magnitude of mechanism replacement can produce''---this is a
local sensitivity analysis tool for side-channel defense design.

\subsection{Statistical Manifold and Fisher Information Matrix}

Given a parametric distribution family
$\{p(y;\theta) : \theta\in\Theta\}$, where
$\theta=(\theta^1,\ldots,\theta^d)$ is the parameter vector and
$\Theta\subset\R^d$ is the parameter space, this family naturally
constitutes a $d$-dimensional differentiable manifold, called the
\emph{statistical manifold}.

In the side-channel analysis context, the parameter $\theta$ should
correspond to observable traffic statistics.
For example, for a packet-length distribution (e.g., exponential or Gamma),
$\theta$ may take the mean parameter $\mu_{\ell}$ and shape parameter
$k_{\ell}$; for a packet inter-arrival distribution (e.g., a Poisson
process), $\theta$ may take the arrival rate $\lambda$; for defense
mechanisms, $\theta$ may include the padding rate $\alpha$ and timing
perturbation scale $\tau_{\mathrm{time}}$.

\begin{definition}[Fisher Information Matrix]\label{def:fim}
For a distribution $p(y;\theta)$, the Fisher information matrix
$G(\theta)=[g_{ij}(\theta)]$ is defined as
\begin{equation}\label{eq:fim}
g_{ij}(\theta) = \E_{Y\sim P_\theta}\!\left[\frac{\partial\log p(Y;\theta)}{\partial\theta^i}
\frac{\partial\log p(Y;\theta)}{\partial\theta^j}\right].
\end{equation}
\end{definition}

The Fisher information matrix gives the Riemannian metric of the
statistical manifold
$\mathcal{M}=\{p(y;\theta):\theta\in\Theta\}$:
for tangent vectors $v,w\in T_\theta\mathcal{M}$, their inner product is
$\langle v,w\rangle_\theta = v^T G(\theta) w$.
The geometric meaning is that the larger the eigenvalue of $G(\theta)$ in
the direction $\theta$, the greater the statistical distance caused by a
small change in parameter $\theta$.
In the side-channel context, $G(\theta)$ is a \emph{distinguishability
sensitivity meter}: if $v^\top G(\theta)v$ is large in some direction $v$,
then a small parameter shift of $\delta$ along $v$ produces a significant
distributional difference (KL divergence $\sim\frac{1}{2}\delta^2 v^\top G(\theta)v$),
and the adversary needs only a limited number of samples to distinguish;
conversely, directions where $v^\top G(\theta)v$ is small are ``blind
zones'' for the adversary.
Each eigenvalue of $G(\theta)$ corresponds to the signal-to-noise ratio
of a ``side-channel distinguishing channel.''

\subsection{Local Regularity and Geometric Approximation of Mutual Information}

\begin{assumption}[Local Second-Order Regularity]\label{ass:local-euclidean}
The log-likelihood $\ell(y;\theta)=\log p(y;\theta)$ is three-times
differentiable in $\theta$ within a neighborhood of $\theta$, with bounded
third-order derivatives, and the Fisher information matrix $G(\theta)$ is
continuous in that neighborhood.
Under this condition, the KL divergence has a second-order Taylor expansion
in the parameter perturbation $\delta$, with the quadratic term given by
the Fisher information.
\end{assumption}

\begin{proposition}[KL Divergence and Fisher Information]\label{prop:kl-fim}
Under Assumption~\ref{ass:local-euclidean}, for a small perturbation
$\|\delta\|\leq \varepsilon$, the second-order expansion of the KL
divergence satisfies:
\begin{equation}\label{eq:kl-fim}
D_{\mathrm{KL}}\!\big(P_{\theta}\,\|\,P_{\theta+\delta}\big)
  = \frac{1}{2}\delta^\top G(\theta)\delta + O(\|\delta\|^3).
\end{equation}
\end{proposition}

\begin{proof}
From the KL divergence definition, with $\ell(y;\theta)=\log p(y;\theta)$:
\[
D_{\mathrm{KL}}\!\big(P_{\theta}\,\|\,P_{\theta+\delta}\big)
  =\E_{Y\sim P_\theta}\!\left[\ell(Y;\theta)-\ell(Y;\theta+\delta)\right].
\]
Applying a second-order Taylor expansion of $\ell(Y;\theta+\delta)$ at $\theta$:
\[
\ell(Y;\theta+\delta)=\ell(Y;\theta)+\delta^\top \nabla_\theta \ell(Y;\theta)
  +\tfrac{1}{2}\delta^\top \nabla_\theta^2\ell(Y;\theta)\delta+O(\|\delta\|^3).
\]
Taking expectations under $P_\theta$ and using the score zero-mean property
$\E_{P_\theta}[\nabla_\theta \ell(Y;\theta)]=0$ and the Fisher information
identity $G(\theta)=-\E_{P_\theta}[\nabla_\theta^2\ell(Y;\theta)]$:
\[
D_{\mathrm{KL}}\!\big(P_{\theta}\,\|\,P_{\theta+\delta}\big)
  =\tfrac{1}{2}\delta^\top G(\theta)\delta+O(\|\delta\|^3). \qed
\]
\end{proof}

Applying the above local expansion to the conditional distribution
$p(y;\theta_x)$ in a neighborhood of the mixture-distribution parameter
$\bar\theta$, and assuming the conditional distributions belong to the same
regular parametric family with $\sup_x \|\theta_x-\bar\theta\| \leq \varepsilon$
and $G(\bar\theta)$ approximating $G(\theta_x)$ throughout the neighborhood
(error absorbed into $O(\varepsilon^3)$):
\begin{equation}\label{eq:mi-approx}
I(X;Y) = \E_x\!\left[\tfrac{1}{2}(\theta_x-\bar\theta)^\top
  G(\bar\theta)(\theta_x-\bar\theta)\right] + O\!\left(\varepsilon^3\right).
\end{equation}

\begin{remark}
The trace form below is used to characterize the local sensitivity
direction of defense parameters, not as a primary empirical tool.
The mutual information quantities in the empirical section are estimated
directly; the trace approximation is used only for interpretive comparison.
\end{remark}

\begin{proposition}[Covariance-Form MI Approximation]\label{prop:mi-cov}
Define the parameter covariance matrix
$\Sigma_\theta := \E_{x\sim P_X}[(\theta_x-\bar\theta)(\theta_x-\bar\theta)^\top]$.
Under the small-perturbation assumptions above:
\begin{equation}\label{eq:mi-trace}
I(X;Y) = \frac{1}{2\ln 2}\mathrm{Tr}\!\big(G(\bar\theta)\Sigma_\theta\big) + O(\varepsilon^3).
\end{equation}
The factor $\frac{1}{\ln 2}$ converts nats to bits; all mutual
information quantities in this paper are in bits.
\end{proposition}

\begin{proof}
Using the trace property $\mathrm{Tr}(AB) = \mathrm{Tr}(BA)$ and the
interchangeability of trace and expectation:
\begin{align*}
\E_x\!\left[(\theta_x-\bar\theta)^\top G(\bar\theta)(\theta_x-\bar\theta)\right]
  &= \mathrm{Tr}\!\big(G(\bar\theta)\E_x[(\theta_x-\bar\theta)(\theta_x-\bar\theta)^\top]\big)\\
  &= \mathrm{Tr}\!\big(G(\bar\theta)\Sigma_\theta\big).
\end{align*}
Substituting into~\eqref{eq:mi-approx} and dividing by $\ln 2$ gives the
result. \qed
\end{proof}

Equation~\eqref{eq:mi-trace} has a clear geometric physical meaning:
side-channel leakage depends on the product of two factors:
(1) the Fisher information matrix $G(\bar\theta)$, characterizing the local
curvature/sensitivity of the statistical manifold, reflecting the
distinguishability of different directions in parameter space;
(2) the covariance matrix $\Sigma_\theta$, characterizing the spread of
application semantics $X$ in parameter space, reflecting the distinguishability
of different applications by their traffic features.
Leakage is significant only when the high-sensitivity directions of the
Fisher metric coincide with the principal directions of semantic spread.


\begin{proposition}[Fisher Information Chain Rule]\label{prop:fisher-chain}
For observed variables $Y=(Y_{\mathrm{size}},Y_{\mathrm{dir}})$, if the
joint density factorizes as
$p(y_{\mathrm{size}},y_{\mathrm{dir}};\theta)=p(y_{\mathrm{size}};\theta)\,
p(y_{\mathrm{dir}}\mid y_{\mathrm{size}};\theta)$
and both the marginal and conditional distributions are regularly
parametrized by the \emph{same} parameter $\theta$ (satisfying the
usual regularity conditions: support independent of $\theta$, score
function exists and integration and differentiation can be exchanged),
then the Fisher information matrix satisfies:
\begin{equation}\label{eq:fisher-chain}
G_{Y}(\theta) = G_{Y_{\mathrm{size}}}(\theta)
  + \E_{Y_{\mathrm{size}}\sim p(\cdot;\theta)}\!\Big[\,G_{Y_{\mathrm{dir}}\mid Y_{\mathrm{size}}}(\theta;Y_{\mathrm{size}})\,\Big],
\end{equation}
where $G_{Y_{\mathrm{dir}}\mid Y_{\mathrm{size}}}(\theta;y_{\mathrm{size}})$
is the conditional Fisher information matrix given
$Y_{\mathrm{size}}=y_{\mathrm{size}}$.
\end{proposition}

\begin{proof}
The Fisher matrix is
$G_Y(\theta) = \E[\nabla_\theta \log p(Y;\theta)
  \nabla_\theta \log p(Y;\theta)^\top]$.
By the chain rule:
$\nabla_\theta \log p(Y;\theta) = \nabla_{\mathrm{size}} + \nabla_{\mathrm{dir}|\mathrm{size}}$,
where $\nabla_{\mathrm{size}} = \nabla_\theta \log p(Y_{\mathrm{size}};\theta)$
and $\nabla_{\mathrm{dir}|\mathrm{size}} = \nabla_\theta \log p(Y_{\mathrm{dir}}\mid Y_{\mathrm{size}};\theta)$.
Expanding and using $\E[\nabla_{\mathrm{dir}|\mathrm{size}}\mid Y_{\mathrm{size}}]=0$
(score zero-mean for the conditional), the cross terms vanish.
The first term is $G_{Y_{\mathrm{size}}}(\theta)$.
The second term equals
$\E_{Y_{\mathrm{size}}}[G_{Y_{\mathrm{dir}}\mid Y_{\mathrm{size}}}(\theta;Y_{\mathrm{size}})]$
by iterated expectation, yielding the result. \qed
\end{proof}

\end{document}
